\chapter{Da detecção à segmentação}
\label{cap:segmentacao}

\begin{objetivos}
Ao final deste capítulo, você será capaz de:
\begin{itemize}[leftmargin=*,itemsep=0.2em]
  \item distinguir detecção de segmentação e justificar por que a segunda é
        operacionalmente superior;
  \item projetar uma U-Net unidimensional para segmentação de séries temporais;
  \item explicar o papel das conexões de atalho na preservação da resolução temporal;
  \item avaliar segmentação com IoU e com SoftED, e interpretar a diferença entre
        as duas famílias de métrica;
  \item diagnosticar desalinhamento de anotação a partir de padrões anômalos nas
        métricas.
\end{itemize}
\end{objetivos}

\section{A pergunta errada}

Todos os capítulos anteriores fizeram a mesma pergunta: \emph{esta janela contém uma
anomalia?} A resposta é um rótulo por janela.

Considere o que isso significa em operação. Com janelas de \SI{900}{\second}, o sistema
informa ``houve um evento em algum ponto destes quinze minutos''. O operador precisa
então examinar novecentas amostras para descobrir onde. Se a janela for de
\SI{7500}{\second}, são mais de duas horas de sinal.

A pergunta correta é outra: \emph{quais amostras, exatamente, pertencem ao evento?}

\begin{definicao}{Segmentação de séries temporais}
Tarefa de atribuir um rótulo a \emph{cada amostra} de uma série temporal, produzindo uma
saída de mesmo comprimento que a entrada. Formalmente, dado
$\mat{X} \in \R^{T \times S}$, produzir $\vect{y} \in \{0,1\}^{T}$, em contraste com a
detecção, que produz um único $y \in \{0,1\}$.
\end{definicao}

\begin{destaque}[Três ganhos, não um]
A mudança de detecção para segmentação parece um refinamento de granularidade. Ela é
mais que isso.

Primeiro, \textbf{localização}: o operador sabe onde olhar.

Segundo, \textbf{início do evento}: com o instante inicial identificado, torna-se
possível medir antecedência de detecção --- métrica que o rótulo por janela não permite
sequer definir.

Terceiro, e menos óbvio, \textbf{qualidade da entrada para as etapas seguintes}. Um
segmento delimitado com precisão contém apenas o evento; uma janela contém o evento mais
o regime normal que o precede. Todas as estat\'isticas calculadas sobre o segmento ---
e o \cref{cap:agente} calculará treze delas por sensor --- são tão boas quanto a
delimitação. \textbf{A segmentação é pré-requisito da explicabilidade.}
\end{destaque}

\section{A U-Net unidimensional}

\subsection{Por que esta arquitetura}

A U-Net foi proposta por \citet{Ronneberger2015} para segmentação de imagens médicas,
sob restrições que ecoam as deste problema: poucos exemplos anotados, necessidade de
saída na mesma resolução da entrada, e classes desbalanceadas.

Sua ideia central é o contorno de um dilema. Para reconhecer um evento é preciso
contexto amplo --- e contexto amplo se obtém reduzindo a resolução. Mas para delimitá-lo
com precisão é preciso resolução fina. As \termo{conexões de atalho}{skip connections}
resolvem o conflito transportando as ativações de alta resolução do codificador
diretamente para o decodificador, contornando o gargalo.

\subsection{Arquitetura}

\begin{center}
\small
\begin{tabular}{lll}
\toprule
\textbf{Bloco} & \textbf{Composição} & \textbf{Resolução} \\
\midrule
Entrada     & $(L, 1)$ & $L$ \\
Cod.\ 1     & Conv1D(32,5)+BN, $\times 2$; MaxPool(2); \ing{Dropout} $0{,}2$ & $L/2$ \\
Cod.\ 2     & Conv1D(64,5)+BN, $\times 2$; MaxPool(2); \ing{Dropout} $0{,}2$ & $L/4$ \\
Gargalo     & Conv1D(128,5)+BN, $\times 2$ & $L/4$ \\
Dec.\ 1     & \ing{UpSampling}(2); concatena atalho (64); Conv1D(64,5)+BN, $\times 2$ & $L/2$ \\
Dec.\ 2     & \ing{UpSampling}(2); concatena atalho (32); Conv1D(32,5)+BN, $\times 2$ & $L$ \\
Saída       & Conv1D(1,1), sigmoide & $L$ \\
\bottomrule
\end{tabular}
\end{center}

A saída é uma probabilidade por amostra, limiarizada em $0{,}5$. O modelo é aplicado
\textbf{por sensor}, e as regiões resultantes são combinadas em seguida --- decisão que
preserva a informação de \emph{qual} sensor evidencia o evento, e que será usada pelo
agente do \cref{cap:agente}.

\begin{destaque}[Por que núcleo 5 e não 3]
Com núcleos de tamanho 5, duas convoluções por bloco e dois níveis de subamostragem, o
campo receptivo do gargalo cobre algumas dezenas de amostras --- suficiente para
capturar a dinâmica de transição de um evento de poço, e ainda estreito o bastante para
que a fronteira do segmento não seja borrada. Com núcleos 3, o campo receptivo exigiria
mais níveis; com núcleos 7 ou maiores, as fronteiras perdem definição. É um compromisso
que vale a pena verificar empiricamente em qualquer aplicação nova.
\end{destaque}

\section{Protocolo experimental}

\begin{itemize}[leftmargin=*,itemsep=0.3em]
  \item \textbf{Dados}: 3W complementado por instâncias proprietárias de poços reais;
        quatro variáveis (P-PDG, P-TPT, T-TPT e P-MON-CKP).
  \item \textbf{Amostra}: 200 passos $\times$ 4 variáveis, ou 800 dimensões quando
        achatada.
  \item \textbf{Janela deslizante}: comprimento de 100 passos, passo de \SI{5}{\second}.
  \item \textbf{Pré-processamento}: imputação de faltantes (último válido, próximo
        válido, média), padronização e remoção de ruído por transformada
        \ing{wavelet}.
  \item \textbf{Validação}: cruzada estratificada em cinco dobras.
  \item \textbf{Treinamento}: Adam, 30 a 50 épocas.
  \item \textbf{Limiares}: $0{,}5$ na saída da U-Net; percentil 95 para a distância de
        Mahalanobis (\cref{cap:mahalanobis}); $0{,}3$ para os classificadores binários.
\end{itemize}

A avaliação é feita \textbf{poço a poço}, e não sobre um conjunto agregado. Dez poços
--- identificados como G1, A1, A2, A3, B1, C1, C2, D1, E1 e F1 --- são reportados
individualmente. Essa escolha é deliberada: o desvio padrão entre poços é a única
medida honesta de generalização disponível, e agregá-los produziria um número único que
esconderia exatamente a informação mais relevante.

\section{Resultados}

\subsection{Autoencoder \emph{versus} U-Net}

A \cref{tab:seg-comparacao} confronta a abordagem anterior --- detecção por erro de
reconstrução de autoencoder, como no \cref{cap:dual} --- com a segmentação por U-Net,
poço a poço.

\begin{table}[htbp]
\centering
\caption{Detecção por autoencoder \emph{versus} segmentação por U-Net, avaliadas no
nível de amostra, por poço.}
\label{tab:seg-comparacao}
\small
\begin{tabular}{lcccc}
\toprule
& \multicolumn{2}{c}{\textbf{Autoencoder}} & \multicolumn{2}{c}{\textbf{U-Net}} \\
\cmidrule(lr){2-3}\cmidrule(lr){4-5}
\textbf{Poço} & \textbf{ACC} & $\bm{F_1}$ & \textbf{ACC} & $\bm{F_1}$ \\
\midrule
G1 & $0{,}818$ & $0{,}706$ & $0{,}882$ & $0{,}731$ \\
A1 & $0{,}541$ & $0{,}703$ & $0{,}775$ & $0{,}739$ \\
A2 & $0{,}582$ & $0{,}736$ & $0{,}929$ & $0{,}936$ \\
A3 & $0{,}293$ & $0{,}454$ & $0{,}951$ & $0{,}915$ \\
B1 & $0{,}285$ & $0{,}443$ & $0{,}752$ & $0{,}632$ \\
C1 & $0{,}772$ & $0{,}766$ & $0{,}885$ & $0{,}817$ \\
C2 & $0{,}329$ & $0{,}277$ & $0{,}975$ & $0{,}912$ \\
D1 & $0{,}296$ & $0{,}457$ & $0{,}892$ & $0{,}777$ \\
E1 & $0{,}496$ & $0{,}663$ & $0{,}870$ & $0{,}874$ \\
F1 & $0{,}567$ & $0{,}724$ & $0{,}720$ & $0{,}674$ \\
\midrule
\textbf{Global} & $0{,}500 \pm 0{,}187$ & $0{,}593 \pm 0{,}160$ &
                  $\bm{0{,}863 \pm 0{,}082}$ & $\bm{0{,}801 \pm 0{,}102}$ \\
\bottomrule
\end{tabular}
\end{table}

\begin{resultadochave}
A acurácia média sobe de $0{,}500$ para $0{,}863$ --- e o desvio padrão entre poços cai
de $0{,}187$ para $0{,}082$, uma redução de $56\,\%$.
\end{resultadochave}

\subsection{O número que importa não é a média}

Uma acurácia de $0{,}500$ em um problema binário é o desempenho de uma moeda. É tentador
concluir que o autoencoder ``não funciona'', mas essa leitura é errada e vale
desmontá-la: o autoencoder do \cref{cap:dual} obteve $0{,}9894$ em operação real.

A diferença está na \emph{tarefa}. Lá ele decidia se uma janela era anômala; aqui, cada
amostra. O erro de reconstrução é um escalar por janela e não carrega informação
temporal fina; usá-lo para segmentar é usá-lo para o que ele não faz.

\begin{destaque}[A redução da variância é o resultado principal]
A média subiu $73\,\%$, mas é o desvio padrão que conta a história decisiva.

Com o autoencoder, a acurácia varia de $0{,}285$ (B1) a $0{,}818$ (G1) --- uma faixa de
cinquenta e três pontos percentuais. O sistema funciona bem em alguns poços e falha em
outros, e não há como saber de antemão em quais. \textbf{Um sistema imprevisível é
inutilizável, mesmo que sua média seja aceitável}, porque o operador não pode calibrar
quanta confiança depositar no alerta.

Com a U-Net, a faixa é de $0{,}720$ (F1) a $0{,}975$ (C2). O pior caso passa a ser
melhor que o caso médio anterior. É essa propriedade --- e não a média --- que torna o
sistema implantável.
\end{destaque}

Vale ainda notar o poço A3: acurácia de $0{,}293$ com o autoencoder e $0{,}951$ com a
U-Net, um ganho de sessenta e seis pontos percentuais. Esse poço reaparecerá adiante,
com um comportamento que revela algo sobre os dados, não sobre os modelos.

\subsection{Qualidade da delimitação}

Acurácia por amostra é uma métrica generosa em problemas desbalanceados. A
\cref{tab:seg-iou} acrescenta a interseção sobre união (\cref{eq:iou}), que mede
diretamente a sobreposição entre o segmento predito e o real.

\begin{table}[htbp]
\centering
\caption{Acurácia e interseção sobre união da segmentação por U-Net, por poço.}
\label{tab:seg-iou}
\small
\begin{tabular}{lcc}
\toprule
\textbf{Poço} & \textbf{ACC (\%)} & \textbf{IoU (\%)} \\
\midrule
A1 & $89{,}00$ & $74$ \\
A2 & $95{,}50$ & $93$ \\
A3 & $96{,}00$ & $85$ \\
B1 & $84{,}00$ & $56$ \\
C1 & $99{,}50$ & $98$ \\
C2 & $97{,}75$ & $84$ \\
D1 & $97{,}05$ & $93$ \\
E1 & $94{,}22$ & $83$ \\
F1 & $86{,}94$ & $72$ \\
G1 & $98{,}73$ & $94$ \\
\midrule
\textbf{Global} & $\bm{94{,}00}$ & $\bm{85}$ \\
\bottomrule
\end{tabular}
\end{table}

\begin{resultadochave}
IoU global de $85\,\%$: o segmento predito e o segmento real coincidem em cerca de cinco
sextos de sua união.
\end{resultadochave}

O poço B1 merece atenção: $84\,\%$ de acurácia mas apenas $56\,\%$ de IoU. A combinação é
diagnóstica --- o modelo acerta a maior parte das amostras, mas as fronteiras do
segmento estão substancialmente deslocadas. Em termos operacionais, o evento é detectado
e mal delimitado. Como toda a cadeia de explicabilidade depende da delimitação, esse é
um caso em que a saída deve ser marcada como de baixa confiança.

\section{Avaliação por SoftED}

\subsection{O problema das métricas ponto a ponto}

Considere um evento que começa na amostra 1000. Um detector o sinaliza na amostra 995;
outro, na 1400. Pela acurácia por amostra, ambos erram --- o primeiro erra cinco
amostras, o segundo quatrocentas. A diferença entre um sistema que antecipa em cinco
segundos e um que atrasa em quase sete minutos praticamente desaparece na métrica.

Em detecção de eventos, isso é inaceitável, e o SoftED de \citet{Salles2024softed}
existe para corrigi-lo, atribuindo crédito parcial a detecções próximas do instante
verdadeiro dentro de uma tolerância $k$.

\subsection{Resultados}

\begin{table}[htbp]
\centering
\caption{Avaliação SoftED com tolerância $k = 1000$ amostras, por poço.}
\label{tab:seg-softed}
\small
\begin{tabular}{lccc}
\toprule
\textbf{Poço} & \textbf{Precisão} & \textbf{Revocação} & $\bm{F_1}$ \\
\midrule
G1 & $1{,}00$ & $1{,}00$ & $1{,}00$ \\
A1 & $1{,}00$ & $1{,}00$ & $1{,}00$ \\
C2 & $1{,}00$ & $1{,}00$ & $1{,}00$ \\
D1 & $1{,}00$ & $1{,}00$ & $1{,}00$ \\
F1 & $1{,}00$ & $1{,}00$ & $1{,}00$ \\
B1 & $1{,}00$ & $0{,}40$ & $0{,}58$ \\
E1 & $1{,}00$ & $0{,}41$ & $0{,}58$ \\
A2 & $1{,}00$ & $0{,}30$ & $0{,}46$ \\
A3 & $0{,}00$ & $1{,}00$ & $0{,}00$ \\
\midrule
\textbf{Global} & $\bm{0{,}89}$ & $\bm{0{,}68}$ & $\bm{0{,}74}$ \\
\bottomrule
\end{tabular}
\end{table}

Métricas complementares globais: antecedência de detecção de $0{,}22$ e taxa de falso
positivo de $0{,}11$.

\subsection{Três leituras}

\paragraph{Precisão alta, revocação baixa.} Precisão de $0{,}89$ contra revocação de
$0{,}68$: quando o sistema aponta um evento, ele quase sempre está certo; mas ele deixa
passar cerca de um terço dos eventos. Para um sistema de alerta operacional, essa é a
assimetria preferível --- alarmes falsos destroem a confiança do operador mais
rapidamente do que eventos perdidos, que ainda podem ser capturados por outros meios.
Mas $0{,}68$ é um número que precisa melhorar.

\paragraph{Cinco poços perfeitos.} G1, A1, C2, D1 e F1 obtêm $1{,}00$ nas três métricas.
Metade dos poços é resolvida completamente. Isso indica que o método não tem uma
limitação intrínseca --- ele falha por razões específicas, em poços específicos, e essas
razões podem ser investigadas.

\paragraph{Antecedência de $0{,}22$.} O sistema raramente antecipa o evento; ele o
detecta no momento em que ocorre ou pouco depois. Para eventos de desenvolvimento lento
--- incrustação, aumento de BSW --- a antecipação seria valiosa e não está sendo obtida.

\subsection{O poço A3, ou como uma métrica revela um problema de dados}

O poço A3 apresenta a combinação mais estranha do livro:

\begin{center}
\small
\begin{tabular}{lc}
\toprule
Acurácia de segmentação (\cref{tab:seg-comparacao}) & $0{,}951$ \\
IoU (\cref{tab:seg-iou})                            & $85\,\%$ \\
SoftED --- revocação                                 & $1{,}00$ \\
SoftED --- precisão                                  & $\bm{0{,}00}$ \\
\bottomrule
\end{tabular}
\end{center}

Um modelo que segmenta com $95\,\%$ de acurácia e $85\,\%$ de IoU, e cuja precisão de
detecção de evento é exatamente zero. Isso não é um resultado ruim --- é um resultado
\emph{impossível}, sob a hipótese de que as anotações estão corretas.

\begin{destaque}[O que precisão $0{,}00$ com revocação $1{,}00$ significa]
Revocação $1{,}00$ indica que todo evento anotado tem uma detecção próxima. Precisão
$0{,}00$ indica que nenhuma detecção está próxima de um evento anotado. As duas
afirmações só são simultaneamente verdadeiras se houver \textbf{muitas} detecções, das
quais apenas algumas caem perto de eventos anotados --- ou se houver um deslocamento
sistemático maior que a tolerância entre o que o modelo detecta e o que está anotado.

A segunda hipótese é a mais provável, e ela aponta para \emph{desalinhamento da
anotação} do poço A3, não para falha do modelo. A alta acurácia de segmentação
corrobora: o modelo encontra consistentemente uma estrutura no sinal --- ela apenas não
está onde o rótulo diz.
\end{destaque}

Este é o tipo de achado que justifica reportar múltiplas métricas em vez de uma. Nenhuma
delas, isoladamente, revelaria o problema; a \emph{contradição} entre elas revela.

\section{Os classificadores binários sobre segmentos}

Com o segmento delimitado, os classificadores binários do \cref{cap:binarios} podem ser
aplicados a ele, e não à janela inteira. A arquitetura, mais leve por operar sobre um
segmento já isolado, é:

\begin{figure}[htbp]
\centering
\includegraphics[width=\textwidth]{Figure_4.png}
\caption{Classificador binário aplicado ao segmento identificado, com as dimensões
efetivas de cada tensor de parâmetros: uma camada recorrente seguida de camadas densas
decrescentes até uma única saída escalar. A figura mostra a variante com LSTM de 256
unidades; o número de unidades recorrentes é ajustado por classe. Adaptado de
\citet{Lopes2026OTC}.}
\label{fig:seg-classificador}
\end{figure}

\begin{center}
\small
\begin{tabular}{ll}
\toprule
\textbf{Camada} & \textbf{Configuração} \\
\midrule
Entrada       & segmento \\
LSTM(128 a 256) & opcional, conforme a classe \\
Densa(64)     & + BN + \ing{Dropout} $0{,}3$ \\
Densa(32)     & + BN + \ing{Dropout} $0{,}2$ \\
Densa(16)     & + \ing{Dropout} $0{,}1$ \\
Densa(1)      & sigmoide \\
\bottomrule
\end{tabular}
\end{center}

O desempenho desses classificadores, e sobretudo o efeito do \emph{espaço} em que o
segmento é representado, é o assunto do \cref{cap:mahalanobis}.

\section{Limitações}

\begin{enumerate}[leftmargin=*,itemsep=0.4em]
  \item \textbf{Revocação de $0{,}68$}. Um terço dos eventos não é detectado dentro da
        tolerância.
  \item \textbf{Antecedência de $0{,}22$}. Praticamente não há detecção antecipada.
  \item \textbf{Variabilidade entre poços persistente}, ainda que reduzida: IoU de
        $56\,\%$ em B1 contra $98\,\%$ em C1.
  \item \textbf{Dependência da qualidade da anotação}, como o caso A3 demonstra de forma
        contundente.
  \item \textbf{Segmentação binária}. A U-Net separa normal de anômalo; ela não
        classifica o tipo de anomalia.
\end{enumerate}

\resumocapitulo{%
Trocar detecção por segmentação --- rótulo por amostra em vez de rótulo por janela ---
produz três ganhos: localização do evento, identificação do instante inicial e, o mais
importante para o restante do livro, qualidade das estatísticas que alimentarão a
camada de explicabilidade. A U-Net unidimensional, com codificador de dois níveis,
gargalo de 128 filtros e conexões de atalho que preservam a resolução temporal, eleva a
acurácia média de $0{,}500$ para $0{,}863$ e reduz o desvio padrão entre poços de
$0{,}187$ para $0{,}082$ --- é a redução de variância, mais que o ganho de média, que
torna o sistema implantável. A interseção sobre união global é de $85\,\%$. A avaliação
por SoftED, que atribui crédito parcial a detecções próximas do instante verdadeiro,
dá precisão de $0{,}89$, revocação de $0{,}68$ e medida $F_1$ de $0{,}74$, com cinco dos
dez poços resolvidos perfeitamente. O poço A3, com revocação $1{,}00$ e precisão
$0{,}00$ apesar de $95\,\%$ de acurácia de segmentação, revela desalinhamento de
anotação --- um achado que só a contradição entre métricas distintas torna visível.}

\begin{exercicios}
\item Calcule o campo receptivo do gargalo da U-Net descrita, considerando núcleos de
      tamanho 5, duas convoluções por bloco e dois níveis de subamostragem por fator 2.
      Compare com a duração típica da transição de um fechamento de DHSV.

\item Implemente a U-Net e reproduza a \cref{tab:seg-comparacao}. Em seguida, remova as
      conexões de atalho e repita. Quantifique a degradação em acurácia e, separadamente,
      em IoU. Qual das duas métricas é mais sensível à remoção? Por quê?

\item Um modelo obtém acurácia de $0{,}84$ e IoU de $0{,}56$ (poço B1). Construa um
      exemplo de segmentação predita e real, sobre uma série de mil amostras, compatível
      com esse par de valores. O que isso revela sobre o erro do modelo?

\item Implemente o SoftED e avalie a sensibilidade dos resultados à tolerância $k$. Trace
      precisão, revocação e $F_1$ globais em função de $k$, de 100 a 5000. Existe um
      valor de $k$ além do qual os resultados deixam de mudar? O que ele significa
      fisicamente?

\item O poço A3 apresenta precisão $0{,}00$ e revocação $1{,}00$. Projete um procedimento
      automático que detecte esse padrão de contradição entre métricas e sinalize
      possíveis erros de anotação em um conjunto de dados novo.

\item A antecedência de detecção é de $0{,}22$. Proponha três modificações da arquitetura
      ou do treinamento que poderiam aumentá-la --- por exemplo, ponderação assimétrica
      da perda nas amostras iniciais do evento --- e discuta o efeito esperado de cada
      uma sobre a taxa de falso positivo.

\item A U-Net é aplicada por sensor e as regiões são combinadas depois. Compare essa
      estratégia com uma U-Net multivariada que recebe os quatro sensores
      simultaneamente. Que informação cada abordagem preserva e qual descarta? Implemente
      ambas e compare.
\end{exercicios}

\begin{leituras}
\item \citet{Lopes2026OTC} --- o artigo que originou este capítulo.
\item \citet{Ronneberger2015} --- a U-Net original, cuja motivação em imagens médicas
      espelha as restrições do problema de poços.
\item \citet{Salles2024softed} --- a métrica SoftED e a discussão sobre por que métricas
      ponto a ponto são inadequadas para detecção de eventos.
\item \citet{VazVargas2026_3w2} --- a versão 2.0 do 3W, com discussão sobre a qualidade das
      anotações.
\end{leituras}
